\subsubsection{Ablation on Top-$k$ Exponent Coding}
\label{sec:ablation_topk}

\begin{wraptable}{r}{0.35\linewidth}
    \vspace{-1.2em}
    \centering
    \caption{Top-$k$ exponent coding.}
    \label{tab:ablation_topk}
    \footnotesize
    \setlength{\tabcolsep}{4pt}
    \renewcommand{\arraystretch}{1.08}
    \begin{tabular}{lcc}
        \toprule
        \textbf{Metric} & \textbf{Top-8} & \textbf{Top-16} \\
        \midrule
        Code width & 3-bit & 4-bit \\
        Coverage & 95.62\% & 99.67\% \\
        Ratio & $1.364\times$ & $1.316\times$ \\
        Encode(GB/s) & 272.6 & 329.9 \\
        Decode(GB/s) & 492.2 & 1263.3 \\
        Escape rate & 4.38\% & 0.33\% \\
        \bottomrule
    \end{tabular}
    \vspace{-1.7em}
\end{wraptable}

Table~\ref{tab:ablation_topk} compares Top-8 3-bit coding with Top-16 4-bit coding.
Top-8 slightly improves the compression ratio from $1.316\times$ to $1.364\times$, but substantially reduces throughput.
Top-16 is $1.21\times$ faster in encoding and $2.57\times$ faster in decoding.

The slowdown of Top-8 comes from both higher escape overhead and less GPU-friendly packing.
Reducing the codebook from 16 to 8 entries lowers coverage from 99.67\% to 95.62\%, increasing the escape rate from 0.33\% to 4.38\%.
This nearly $13\times$ higher escape rate increases memory-intensive escape collection and sparse overwrite costs.
In addition, 3-bit codes do not align naturally with byte-oriented GPU operations, unlike 4-bit codes where two exponents form one byte.
The FP8 results in Table~\ref{tab:fp8_exact_results} show the same trend: lower-bit exponent coding does not necessarily improve end-to-end performance.